\subsection{The Lexicon and Lexical Relations}
\label{sec:lexical-relations}
\slidesrc{Semantics (long version): 11--21}

\subsubsection{Lexicon and lexical units}
\slidesrc{Semantics (long version): 11--13}

Linguistic theories generally regard a human language as made of two parts:
\begin{itemize}[nosep]
  \item a \textbf{lexicon}: essentially a catalogue of the language's words (its
        \emph{wordstock}). It is also thought to include bound morphemes, which cannot
        stand alone as words (most affixes, Section~\ref{sec:morphology}), and in some
        analyses compound words and certain idioms and collocations;
  \item a \textbf{grammar}: a system of rules that combine those words into meaningful
        sentences (Section~\ref{sec:grammars-parsing}).
\end{itemize}
Dictionaries try to list the lexicon in alphabetical order, usually without the bound
morphemes.

\begin{nbdef}{Lexical semantics}
The units of analysis of \textbf{lexical semantics} are \textbf{lexical units}: not only
words but also sub-words (affixes) and larger units (compound words, phrases). Lexical
units make up the lexicon. Lexical semantics studies how the meaning of lexical units
correlates with the structure of the language (syntax), and how meanings relate to each
other through \textbf{lexical relations}: hyponymy, hypernymy, synonymy, antonymy,
homonymy, \ldots
\end{nbdef}

\subsubsection{Synonymy}
\slidesrc{Semantics (long version): 14}

\begin{nbdef}{Synonymy}
\textbf{Synonymy} is similarity of meaning. Leibniz: two expressions are synonymous if
substituting one for the other \emph{never} changes the truth value of a sentence in which
the substitution is made.
\end{nbdef}

Such \emph{global} synonymy is rare (a language with perfect synonyms would be
redundant). A more useful notion is \textbf{synonymy relative to a context}: two
expressions are synonymous in a linguistic context $C$ if substituting one for the other
in $C$ does not alter the truth value. E.g.\ \emph{board} and \emph{plank} are
synonyms in ``he nailed a \underline{\quad} to the wall'' but not in ``the
\underline{\quad} of directors''. A consequence of defining synonymy by substitution:
words of different syntactic categories cannot be synonyms (\emph{quick} and
\emph{quickly} cannot replace each other).

\subsubsection{Antonymy}
\slidesrc{Semantics (long version): 15}

The antonym of a word $x$ is sometimes \emph{not-$x$}, but not always: \emph{rich} and
\emph{poor} are antonyms, but \emph{not rich} does not imply \emph{poor} (many people
consider themselves neither). Antonymy is a lexical relation between \emph{word forms},
not a semantic relation between concepts: \emph{rise}/\emph{fall} and
\emph{ascend}/\emph{descend} are antonym pairs, but \emph{rise}/\emph{descend} sounds
odd, even though the concepts are the same.

\subsubsection{Hyponymy and hypernymy}
\slidesrc{Semantics (long version): 16--17}

\begin{nbdef}{Hyponymy and hypernymy}
\textbf{Hyponymy} is a semantic relation between word meanings: \{maple\} is a
\textbf{hyponym} of \{tree\}. The inverse relation is \textbf{hypernymy}: \{tree\} is a
\textbf{hypernym} of \{maple\}. Also called subordination/superordination,
subset/superset, or the \textbf{IS-A} relation. Test: a native speaker must accept
sentences of the form ``An $x$ is a (kind of) $y$''. For verbs the relation is called
\textbf{troponymy} (\emph{to lisp} is a manner of \emph{to talk}).
\end{nbdef}

\begin{center}
\begin{tikzpicture}[w/.style={font=\small, text=nbblue}, e/.style={draw=nborange!80}]
  \node[w] (an) at (-1.2,0) {animals};
  \node[w] (do) at (-3.8,-1.1) {dogs};
  \node[w] (ca) at (0.4,-1.1) {cats};
  \node[w] (ho) at (2.0,-1.1) {horses};
  \node[w] (pi) at (3.5,-1.1) {pigs};
  \node[w] (te) at (-5.9,-2.2) {terrier};
  \node[w] (bl) at (-3.6,-2.2) {bloodhound};
  \node[w] (co) at (-1.0,-2.2) {cocker spaniel};
  \node[w] (t1) at (-8.7,-3.3) {Cesky terrier};
  \node[w] (t2) at (-5.9,-3.3) {Border terrier};
  \node[w] (t3) at (-2.5,-3.3) {Australian terrier};
  \foreach \c in {do,ca,ho,pi} \draw[e] (an.south) -- (\c.north);
  \foreach \c in {te,bl,co} \draw[e] (do.south) -- (\c.north);
  \foreach \c in {t1,t2,t3} \draw[e] (te.south) -- (\c.north);
\end{tikzpicture}
\captionof{figure}{A hyponymy (IS-A) hierarchy: a Border terrier is a terrier, a terrier
is a dog, a dog is an animal. Hyponymy is transitive.}
\end{center}

\subsubsection{Meronymy and holonymy}
\slidesrc{Semantics (long version): 18--20}

\begin{nbdef}{Meronymy and holonymy}
A concept $C_1$ is a \textbf{meronym} of a concept $C_2$ in a language $L$ if native
speakers of $L$ accept sentences like ``A $C_2$ has a $C_1$ (as a part)'' or ``A $C_1$ is
a part of a $C_2$''. The inverse relation is \textbf{holonymy}. It is the
\textbf{HAS-AS-PART} relation, organised in a part hierarchy (e.g.\ \emph{car} has
\emph{engine}, \emph{door}, \emph{wheel}, \emph{tyres}; \emph{engine} has
\emph{cylinder}, \emph{exhaust valve}, \ldots).
\end{nbdef}

Part-of is asymmetric and, \emph{with caution}, transitive. Transitivity fails when
different kinds of part-whole relation are chained:
\begin{itemize}[nosep]
  \item A musician has an arm. \hfill (component of a body)
  \item An orchestra has a musician. \hfill (member of a group)
  \item but: ?\,An orchestra has an arm.
\end{itemize}

\subsubsection{Homonymy}
\slidesrc{Semantics (long version): 21}

\begin{nbdef}{Homonyms, homophones, homographs}
\begin{itemize}[nosep]
  \item \textbf{Homonyms}: the same word form with multiple unrelated meanings.
  \item \textbf{Homophones}: words that \emph{sound} alike but are written differently.
  \item \textbf{Homographs}: words with the same spelling, possibly different
        pronunciation, and different meanings.
\end{itemize}
\end{nbdef}

\begin{center}
\footnotesize
\begin{tabularx}{\linewidth}{@{}*{3}{>{\raggedright\arraybackslash}X}@{}}
\toprule
\textbf{Homonyms} & \textbf{Homophones} & \textbf{Homographs} \\
\midrule
the \emph{spruce} tree / to \emph{spruce} up & \emph{addition} for maths / \emph{edition} of a book & \emph{desert} = abandon / area of land \\
\emph{suit} yourself / wore a \emph{suit} & I want \emph{to} go / I like it \emph{too} / one plus one is \emph{two} & \emph{bass} = fish / instrument \\
weigh on the \emph{scale} / \emph{scale} the wall & \emph{capitol} building / state \emph{capital} & \emph{close} = nearby / to shut \\
the price is \emph{fair} / go to the \emph{fair} & pick a \emph{flower} / bake with \emph{flour} & \emph{bow} = to bend down / ribbon \\
\bottomrule
\end{tabularx}
\end{center}

\begin{nbwhy}
Homographs and homonyms are a problem for every task that looks at written forms: a
search for \emph{bass} returns both fish and guitars (low precision, Section~\ref{sec:precision-recall}),
and a distributional vector for \emph{bass} mixes the contexts of all its meanings.
Choosing the right meaning in context is word sense disambiguation
(Section~\ref{sec:word-sense-disambiguation}).
\end{nbwhy}
